\subsubsection{Matriz de cossenos diretores}
\label{sec:matriz_cosseno_diretor}

Considerando dois referenciais $A$ e $B$ em rotação relativa. A orientação de $B$ em relação a $A$ pode ser descrita pela matriz de cossenos diretores (\gls{dcm}), denotada por $\mathbf{C}^{B/A}$.
Seguindo a notação de \cite{wie2008space}, cada linha de $\mathbf{C}^{B/A}$ contém as projeções de um versor de $B$ nos eixos de $A$, de modo que

\begin{equation}
\begin{bmatrix}
\vec{b}_1 \\
\vec{b}_2 \\
\vec{b}_3
\end{bmatrix}
=
\mathbf{C}^{B/A}
\begin{bmatrix}
\vec{a}_1 \\
\vec{a}_2 \\
\vec{a}_3
\end{bmatrix},
\label{eq:def_dcm_ba}
\end{equation}

onde $\{\vec{a}_1,\vec{a}_2,\vec{a}_3\}$ e $\{\vec{b}_1,\vec{b}_2,\vec{b}_3\}$ são bases ortonormais de $A$ e $B$, respectivamente. Cada elemento $C_{ij}^{B/A}$ é o cosseno do ângulo entre o $i$-ésimo versor de $B$ e o
$j$-ésimo versor de $A$

\begin{equation}
C_{ij}^{B/A} = \vec{b}_i \cdot \vec{a}_j.
\end{equation}

A matriz $\mathbf{C}^{B/A}$ é, portanto, uma matriz ortogonal no sentido algébrico, suas linhas e colunas são mutuamente ortogonais.
Quando, além disso, cada linha e cada coluna tem norma unitária, diz-se que a matriz é ortonormal.
Neste caso, vale

\begin{equation}
\mathbf{C}^{B/A} \big(\mathbf{C}^{B/A}\big)^{T}
=
\big(\mathbf{C}^{B/A}\big)^{T}\mathbf{C}^{B/A}
=
\mathbf{I},
\label{eq:ortonormalidade_C}
\end{equation}

com $\det \mathbf{C}^{B/A} = +1$ para rotações físicas próprias. Em muitas referências de dinâmica, o termo “matriz ortogonal” já pressupõe esta condição de ortonormalidade.
A relação inversa de \eqref{eq:def_dcm_ba} é escrita como

\begin{equation}
\begin{bmatrix}
\vec{a}_1 \\
\vec{a}_2 \\
\vec{a}_3
\end{bmatrix}
=
\big(\mathbf{C}^{B/A}\big)^{-1}
\begin{bmatrix}
\vec{b}_1 \\
\vec{b}_2 \\
\vec{b}_3
\end{bmatrix}
=
\big(\mathbf{C}^{B/A}\big)^{T}
\begin{bmatrix}
\vec{b}_1 \\
\vec{b}_2 \\
\vec{b}_3
\end{bmatrix},
\label{eq:relacao_inversa_dcm}
\end{equation}

onde foi usada a propriedade ortonormal \eqref{eq:ortonormalidade_C}. Esta matriz não deve ser confundida com a matriz de rotação obtida diretamente a partir de um triplo de ângulos de Euler, para um dado conjunto de ângulos de
Euler existe uma DCM associada, mas a DCM é uma descrição mais geral da orientação.

Segundo \cite{wie2008space}, considera-se que a velocidade angular de $B$ em relação a $A$ é o vetor

\begin{equation}
\vec{\omega}^{B/A}(t)
=
\omega_1(t)\,\vec{b}_1
+
\omega_2(t)\,\vec{b}_2
+
\omega_3(t)\,\vec{b}_3.
\end{equation}

A derivada temporal dos versores de $B$, vista a partir de $A$, é dada pela cinemática de corpo rígido

\begin{equation}
\dot{\vec{b}}_i = \vec{\omega}^{B/A} \times \vec{b}_i,
\qquad i = 1,2,3.
\label{eq:derivada_b_i}
\end{equation}

Partindo de \eqref{eq:relacao_inversa_dcm}, toma-se a derivada temporal em $A$

\begin{equation}
\begin{bmatrix}
\vec{0} \\
\vec{0} \\
\vec{0}
\end{bmatrix}
=
\dot{\mathbf{C}}^{T}
\begin{bmatrix}
\vec{b}_1 \\
\vec{b}_2 \\
\vec{b}_3
\end{bmatrix}
+
\big(\mathbf{C}^{B/A}\big)^{T}
\begin{bmatrix}
\dot{\vec{b}}_1 \\
\dot{\vec{b}}_2 \\
\dot{\vec{b}}_3
\end{bmatrix},
\label{eq:derivada_relacao_inversa}
\end{equation}

onde se usou a notação $\mathbf{C} \equiv \mathbf{C}^{B/A}$ para simplificar.
Substituindo \eqref{eq:derivada_b_i} em \eqref{eq:derivada_relacao_inversa}, obtém-se

\begin{equation}
\dot{\mathbf{C}}^{T}
\begin{bmatrix}
\vec{b}_1 \\
\vec{b}_2 \\
\vec{b}_3
\end{bmatrix}
-
\mathbf{C}^{T}
\begin{bmatrix}
\vec{\omega}\times\vec{b}_1 \\
\vec{\omega}\times\vec{b}_2 \\
\vec{\omega}\times\vec{b}_3
\end{bmatrix}
=
\begin{bmatrix}
\vec{0} \\
\vec{0} \\
\vec{0}
\end{bmatrix}.
\end{equation}

Para manipular o produto vetorial em forma matricial, introduz-se a matriz antissimétrica $\boldsymbol{\Omega}$ associada a $\vec{\omega}$. Esta matriz é o operador que representa o produto vetorial

\begin{equation}
\vec{\omega}\times\vec{v}
=
\boldsymbol{\Omega}\,\vec{v},
\qquad
\boldsymbol{\Omega}
=
\begin{bmatrix}
0      & -\omega_3 &  \omega_2 \\
\omega_3 & 0       & -\omega_1 \\
-\omega_2 & \omega_1 & 0
\end{bmatrix},
\label{eq:def_omega_matriz}
\end{equation}

de modo que $\boldsymbol{\Omega}$ é antissimétrica ($\boldsymbol{\Omega}^{T} = -\boldsymbol{\Omega}$).
Com a definição \eqref{eq:def_omega_matriz}, a expressão anterior pode ser escrita de forma compacta como

\begin{equation}
\big(\dot{\mathbf{C}}^{T} - \mathbf{C}^{T}\boldsymbol{\Omega}\big)
\begin{bmatrix}
\vec{b}_1 \\
\vec{b}_2 \\
\vec{b}_3
\end{bmatrix}
=
\begin{bmatrix}
\vec{0} \\
\vec{0} \\
\vec{0}
\end{bmatrix},
\end{equation}

o que conduz à condição matricial

\begin{equation}
\dot{\mathbf{C}}^{T} - \mathbf{C}^{T}\boldsymbol{\Omega} = \mathbf{0}.
\end{equation}

Tomando a transposta e usando $\boldsymbol{\Omega}^{T} = -\boldsymbol{\Omega}$, obtém-se a equação cinemática diferencial para a matriz de cossenos diretores

\begin{equation}
\dot{\mathbf{C}} + \boldsymbol{\Omega}\,\mathbf{C} = \mathbf{0},
\label{eq:cinematica_dcm}
\end{equation}

conforme apresentado \cite{wie2008space}.
Escrevendo \eqref{eq:cinematica_dcm} elemento a elemento, obtêm-se as nove equações diferenciais em termos dos cossenos diretores $C_{ij}$:

\begin{equation}
\begin{matrix}
    \dot{C}_{11} &= \omega_3 C_{21} - \omega_2 C_{31}, \\
    \dot{C}_{12} &= \omega_3 C_{22} - \omega_2 C_{32}, \\
    \dot{C}_{13} &= \omega_3 C_{23} - \omega_2 C_{33}, \\
    \dot{C}_{21} &= \omega_1 C_{31} - \omega_3 C_{11}, \\
    \dot{C}_{22} &= \omega_1 C_{32} - \omega_3 C_{12}, \\
    \dot{C}_{23} &= \omega_1 C_{33} - \omega_3 C_{13}, \\
    \dot{C}_{31} &= \omega_2 C_{11} - \omega_1 C_{21}, \\
    \dot{C}_{32} &= \omega_2 C_{12} - \omega_1 C_{22}, \\
    \dot{C}_{33} &= \omega_2 C_{13} - \omega_1 C_{23}.
\end{matrix}
\end{equation}

Essas equações mostram explicitamente como cada cosseno diretor evolui em
função da velocidade angular $\vec{\omega}$.
Quando $\vec{\omega}(t)$ é conhecida como função do tempo, a orientação de $B$ em relação a $A$ pode ser obtida pela integração numérica de \eqref{eq:cinematica_dcm}.
A condição de ortonormalidade \eqref{eq:ortonormalidade_C} é um integral de movimento da equação cinemática: se $\mathbf{C}(0)$ é ortonormal, então a solução exata de \eqref{eq:cinematica_dcm} satisfaz $\mathbf{C}(t)\mathbf{C}^{T}(t)=\mathbf{I}$ para todo $t>0$. Na prática, essa propriedade é usada como verificação de consistência numérica e, quando necessário, como critério para re-ortonormalizar a matriz integrada.

Por fim, a comparação entre a DCM $\mathbf{C}^{B/A}$ e a matriz de rotação obtida a partir de ângulos de Euler é fundamental para entender as diferenças entre representações de atitude. A DCM atua diretamente nas bases dos referenciais, enquanto a matriz de rotação em ângulos de Euler é construída como produto de três rotações elementares em torno de eixos prescritos. 